\begin{table}[H]
  \begin{center}
    \setlength\extrarowheight{-6pt}
    \begin{tabular}{c|ccccccccc}
      $n/k$ & 0 & 1 & 2 & 3 & 4 & 5 & 6 & 7 & 8 \\ [3px]
      \hline
      0 & 1 &   &   &   &   &   &   &   &   \\
      1 & 3 & 1 &   &   &   &   &   &   &   \\
      2 & 9 & 6 & 1 &   &   &   &   &   &   \\
      3 & 27 & $\tfrac{109}{4}$ & 9 & 1 &   &   &   &   &   \\
      4 & 81 & 111 & 55 & 12 & 1 &   &   &   &   \\
      5 & 243 & $\tfrac{6841}{16}$ & 285 & $\tfrac{185}{2}$ & 15 & 1 &   &   &   \\
      6 & 729 & $\tfrac{12753}{8}$ & 1351 & 585 & 140 & 18 & 1 &   &   \\
      7 & 2187 & $\tfrac{372709}{64}$ & 6069 & $\tfrac{53011}{16}$ & 1050 & $\tfrac{791}{4}$ & 21 & 1 &   \\
      8 & 6561 & $\tfrac{167943}{8}$ & 26335 & $\tfrac{35049}{2}$ & 6951 & 1722 & 266 & 24 & 1 \\
    \end{tabular}
  \end{center}
  \caption{Values of generalized central factorial numbers $T_3(n,k)$.
    See also the sequences \href{https://oeis.org/A000000}{\texttt{A000000}},
    and \href{https://oeis.org/A000000}{\texttt{A000000}} in the OEIS~\cite{sloane2003line}.
  }
  \label{tab:central-factorial-numbers-t3}
\end{table}
